\section{SoLM：融合 AR 与扩散的排序语言模型}

\subsection{核心思路：排序 + 因果注意力}

SoLM（Sorting-based Language Model，亦称 Esoteric Language Model）的核心贡献是为 MDLM 引入 KV 缓存支持。关键思路：

\begin{importantbox}{SoLM 的训练策略}
将 MDLM 的双向注意力替换为\textbf{因果注意力}，通过以下操作实现：
\begin{enumerate}
    \item 对噪声序列 $\mathbf{z}_t$ 进行\textbf{排序/洗牌}
    \item 将所有\textbf{干净 token 放到左侧}（作为已知上下文）
    \item 将所有\textbf{mask token 放到右侧}（作为待预测目标）
    \item 对排序后的序列施加\textbf{因果注意力}
\end{enumerate}
\end{importantbox}

在这种排列下：
\begin{itemize}
    \item 第一个 mask token 只依赖干净上下文
    \item 第二个 mask token 依赖干净上下文 + 前一个 mask token
    \item 以此类推...
    \item 没有任何 mask token 依赖右侧的"未来" mask token
\end{itemize}

\subsection{推理过程与 KV 缓存}

SoLM 的推理过程充分利用了因果注意力带来的 KV 缓存能力：

\begin{enumerate}
    \item 祖先采样器决定要去噪哪些 mask 位置（例如位置 3 和 2）
    \item \textbf{只对这些 mask token 做前向传播}，附带对应的位置编码
    \item 已预测 token 的激活值被\textbf{冻结}（缓存）
    \item 下一步去噪时，只需处理新的 mask token，利用缓存的 KV
\end{enumerate}

\begin{importantbox}{SoLM 推理的计算节省}
对比 MDLM 每步需要处理整个序列（包括大量不携带信息的 mask token），SoLM 每步\textbf{只处理待去噪的 token}。这带来了两方面的节省：
\begin{enumerate}
    \item 前向传播的序列长度大幅缩短
    \item KV 缓存避免了重复计算已处理 token 的激活
\end{enumerate}
这两个因素共同实现了\textbf{数量级的推理加速}。
\end{importantbox}

\subsection{AR 与扩散的插值}

SoLM 的另一个重要贡献是它能够在纯 AR 模式和纯扩散模式之间\textbf{平滑插值}：

\begin{itemize}
    \item \textbf{纯 AR 模式}：每步只去噪一个 token（从左到右），质量最高但速度最慢
    \item \textbf{纯扩散模式}：每步去噪多个 token（并行），速度最快但可能有质量损失
    \item \textbf{中间模式}：根据需求在速度和质量之间灵活权衡
\end{itemize}

\begin{knowledgebox}{灵活的速度-质量权衡}
这种插值能力意味着用户可以根据具体应用场景的需求来选择最优的工作点。例如，在对延迟敏感的在线服务中可以选择更激进的并行生成，在对质量要求极高的场景中则采用接近 AR 的模式。
\end{knowledgebox}

\subsection{本章小结}

SoLM 通过排序 + 因果注意力的巧妙设计，在保持扩散模型并行生成能力的同时实现了 KV 缓存支持。这使得推理效率实现了数量级的提升，真正使扩散语言模型具备了实用化的推理速度。

\newpage
